\documentclass[conference]{IEEEtran}
\usepackage{times}
\usepackage{amsmath}
\usepackage{booktabs}
\usepackage{balance}
\usepackage[bookmarks=true]{hyperref}

\hypersetup{
  pdftitle={Rebuttal for CoRL Submission 18},
  pdfauthor={},
  pdfsubject={CoRL rebuttal},
  pdfkeywords={}
}
\newcommand{\newparagraph}[1]{{\bf #1}}

\begin{document}

\title{Rebuttal for CoRL Submission \#18\vspace{-1cm}}
\maketitle
\thispagestyle{empty}
\balance

We thank the reviewers and Area Chair for identifying important ambiguities in
the optimization, masking, ablations, and real-robot protocol. We agree with
these concerns and describe concrete corrections and controlled experiments
below. The current results are preliminary sanity checks; final controlled
results will replace the marked entries.

\begin{table}[h]
  \centering
  \caption{SWIPE success/rollouts with 95\% Wilson intervals. ``Pilot'' denotes
  the current hand-stabilized sanity runs; the aligned fixed-mount evaluation
  is ongoing.}
  \label{tab:rebuttal_swipe}
  \begin{tabular}{lcc}
    \toprule
    Policy & Pilot ($N{=}5$) & Fixed mount \\
    \midrule
    Legacy VPCE-DP (reference) & $4/5\ [0.38,0.96]$ & TBD \\
    Clean full, $\alpha{=}2$   & $2/5\ [0.12,0.77]$ & TBD \\
    Clean full, $\alpha{=}1$   & $3/5\ [0.23,0.88]$ & TBD \\
    Clean no-gating            & $0/5\ [0.00,0.43]$ & TBD \\
    Phase-conditioned DP       & $2/5\ [0.12,0.77]$ & TBD \\
    Progress-conditioned DP    & $3/5\ [0.23,0.88]$ & TBD \\
    \bottomrule
  \end{tabular}
\end{table}

The pilot is reported for transparency, not as a model ranking: the first
clean-full batch used a lower execution speed, while the legacy reference used
a different workspace clip. The final column uses one frozen controller,
identical near-contact initialization, and a rigidly mounted reader.

\newparagraph{Q1: Is VPCE only task progress or phase?}
We trained same-backbone Phase-DP and Progress-DP baselines that replace the
VPCE horizon with a progress-derived four-bin phase or normalized trajectory
progress and do not use tactile-derived contact supervision. Their pilot point
estimates are comparable to the clean full models, and all intervals are wide;
therefore these data do not yet establish that VPCE is superior to explicit
progress. We will base the revision on the aligned fixed-mount results. If the
baselines remain competitive, we will narrow the claim: on SWIPE, contact
timing overlaps task phase, while VPCE predicts it online from RGB and motion
without requiring oracle trajectory progress.

\newparagraph{Q2--Q3: Loss weighting and proprioceptive masking.}
We agree that a learned scalar optimized through the same weighted denoising
objective can be driven to its lower bound. We replace it with the fixed,
bounded schedule $w(c)=\operatorname{clip}(1+\alpha c,1,w_{\max})$ and test
$(\alpha,w_{\max})=(2,3)$ and $(1,2)$. We also agree that the original
post-encoder low-dimensional mask was too coarse. The revised policy always
retains proprioception and gripper width and hard-masks only tactile/contact
features. These are direct implementation corrections rather than a defense
of the original choices.

\newparagraph{Q4: Is feature gating necessary?}
The submitted pilot gave $9/10$ for both full and no-gating variants, suggesting
that gating may be auxiliary. The revised no-gating sanity run ($0/5$) points
in the opposite direction, but its small sample and the pilot protocol do not
support a causal conclusion. We therefore will not claim that gating is either
necessary or dispensable until the aligned results are complete, and will move
the component ablations into the main text.

\newparagraph{Human-held props and limitations.}
We agree that hand compliance can assist alignment, even without deliberate
trajectory correction. We therefore separate hand-stabilized pilots from the
rigid-mount evaluation and use the latter to answer this concern. Every SWIPE
trial uses a 60-s timeout and succeeds only when the reader's internal roller
completes one effective rotation. Finally, VPCE is target-adapted using
task-specific tactile demonstrations and calibrated contact labels; we do not
claim zero-shot transfer to unseen tasks, objects, or viewpoints.

\end{document}
